% From mitthesis package
% Version: 1.01, 2023/06/19
% Documentation: https://ctan.org/pkg/mitthesis
%
% The abstract environment creates all the required headers and footnote. 
% You only need to add the text of the abstract itself.
%
% Approximately 500 words or less; try not to use formulas or special characters
% If you don't want an initial indentation, do \noindent at the start of the abstract


% Comments from Martyn

% In general a good project idea, but it needs more evidence, motivation, and support for some of the claims you make.
% Your reference section needs some additional work to ensure it is correct and complete.

% It is not clear what programming language you will be using, and also whether the code is substantial enough since you
% are using so many pre-built/off the shelf solutions. Ensure that when these are removed, that you have still developed
% the majority of the code yourself, and that it is substantial.

% By discussing smart devices in relation specifically to the AirBnB case study, it appears you might have found it
% difficult to find much on the topic.

% You might consider doing some further reading and talking more generally about sustainability and the use of
% technology, smart devices for environmental monitoring before moving on to the detail specific to AirBnB. This will
% help you to find more references for the Critical review that cover some of the research in this area. You can then
% argue that smart devices are a good option for your project due to their benefits, which include cost, small form
% factor, and connectivity. 

This project harnesses the capabilities of smart monitoring devices and cloud computing to systematically 
collect data on energy consumption, water usage, and temperature within an 
AirBnB \cite{Airbnb2024a} property. The data serves to provide guests with 
insights into the environmental impact of their resource use, delivered through daily reports and a comprehensive report 
at the conclusion of their stay. Additionally, the project aimed to develop a web application that enables hosts to 
monitor and manage their property's electricity and water consumption effectively. An 
Application Programming Interface (API) \cite{Amazon2024} has been 
developed to facilitate interaction with the sensor-collected data, allowing both application developers and third-party 
users to have access to datasets for future research and analysis. The overarching goal of this initiative is to expand 
the network of environmentally conscious AirBnB hosts and guests by raising awareness of the environmental implications 
of energy and water consumption during stays.

Make energy and water usage more transparent for AirBnB guests to adjust their behavour while being hosted.